\documentclass[11pt,a4paper]{article}
\usepackage[T1]{fontenc}
\usepackage[utf8]{inputenc}
\usepackage{lmodern,microtype}
\usepackage[a4paper,margin=25mm,headheight=15pt]{geometry}
\usepackage{amsmath,amssymb,amsthm,mathtools}
\usepackage{booktabs,longtable,array,enumitem}
\usepackage[dvipsnames]{xcolor}
\usepackage{fancyhdr}
\usepackage[unicode,colorlinks=true,linkcolor=MidnightBlue,citecolor=MidnightBlue,urlcolor=MidnightBlue,breaklinks=true]{hyperref}
\usepackage{bookmark}
\usepackage{xurl}
\pagestyle{fancy}
\fancyhf{}
\fancyhead[L]{\small Sharp weighted type under formal reversion}
\fancyhead[R]{\small Research article}
\fancyfoot[C]{\thepage}
\renewcommand{\headrulewidth}{0.3pt}
\setlength{\parskip}{3pt}
\setlength{\emergencystretch}{2em}
\setlist{itemsep=3pt,topsep=5pt,leftmargin=1.6em}
\numberwithin{equation}{section}
\newtheorem{theorem}{Theorem}[section]
\newtheorem{lemma}[theorem]{Lemma}
\newtheorem{proposition}[theorem]{Proposition}
\newtheorem{corollary}[theorem]{Corollary}
\newtheorem{inputtheorem}[theorem]{Repository input}
\theoremstyle{definition}
\newtheorem{definition}[theorem]{Definition}
\newtheorem{example}[theorem]{Example}
\theoremstyle{remark}
\newtheorem{remark}[theorem]{Remark}
\newcommand{\C}{\mathbb C}
\newcommand{\R}{\mathbb R}
\newcommand{\N}{\mathbb N}
\newcommand{\typ}{\tau_M}
\newcommand{\Tg}{T_s}
\newcommand{\B}{\mathcal B}
\newcommand{\res}{\operatorname{res}}
\newcommand{\ord}{\operatorname{ord}}
\newcommand{\ee}{\mathrm e}
\newcommand{\dd}{\mathrm d}
\newcommand{\Q}{\widehat Q}
\newcommand{\U}{\widehat U}
\newcommand{\coef}[2]{[#1^{#2}]}
\newcommand{\commit}{\nolinkurl{0c973d8f5e7ef4550f4df1bf486d890037fd9f14}}
\hypersetup{pdftitle={Sharp Weighted Type Under Formal Reversion},pdfauthor={AI-assisted research prepared for Vladimir Reshetnikov},pdfsubject={Exact weighted type, formal inversion, feedback transseries, Borel growth, non-D-finiteness}}
\begin{document}
\begin{titlepage}
\centering
\vspace*{13mm}
{\Large\scshape Transseries and nonlinear inversion\par}
\vspace{12mm}
{\Huge\bfseries Sharp Weighted Type\\[3mm] Under Formal Reversion\par}
\vspace{9mm}
{\Large Zero-loss inversion, exact Borel radii, and\\[2mm]
non-D-finiteness in countable exponential feedback\par}
\vspace{12mm}
{\large Prepared for Vladimir Reshetnikov\par}
\vspace{3mm}
29 September 2026
\vspace{13mm}
\begin{minipage}{0.9\textwidth}
\begin{abstract}
Let $M_0=1$ and let $(M_n)$ be a positive log-convex sequence. We prove
that the weighted coefficient type
\[
 \tau_M(f)=\limsup_{n\to\infty}
       \left(\frac{|[z^n]f|}{M_n}\right)^{1/n}
\]
is preserved by every tangent-to-identity formal compositional inverse
if and only if $M_n^{1/n}\to\infty$. A finite Lagrange majorant gives an
explicit overhead $\exp\{O(n/\sqrt{M_n^{1/n}})\}$, hence no loss in the
exponential constant. Analytic coordinate changes obey an exact scaling
law, without a moderate-growth assumption on the weight.

Using two explicitly stated forward estimates from ProveIt's
countable exponential-feedback report, we obtain the missing sharp type
of its inverse, exact inverse Borel radii, and a logarithmically refined
inverse type. For $\lambda_j=j^2$, this yields
\[
 \limsup_{n\to\infty}
 \left(\frac{|[u^n]\widehat Q(u)|\,\log(n+\ee)^{2n}}{n!}\right)^{1/n}=4.
\]
Its ordinary Borel transform is entire but has infinite classical order,
with an exact iterated-logarithmic maximum-modulus limsup. An elementary
holonomic obstruction proves that both the forward series and its inverse
are non-D-finite. All new arguments are proved in the article; the two
repository inputs, finite computational checks, and remaining research
questions are distinguished explicitly.
\end{abstract}
\end{minipage}
\vfill
\begin{minipage}{0.9\textwidth}\small
\textbf{Status.} Conventional mathematical proofs and reproducible exact
finite checks; not Lean-verified or peer-reviewed. The theorem is a proposed
contribution and resolves a documented repository-level gap. Global
publication priority, and a solution of a named longstanding conjecture,
are not claimed.
\end{minipage}
\end{titlepage}
\setcounter{tocdepth}{1}
\tableofcontents
\vspace{9mm}
\noindent\textbf{Reading guide.}
Sections~\ref{sec:weights}--\ref{sec:weightspecial} form the independent
core: a finite coefficient estimate, exact inversion invariance, its
necessary-and-sufficient weight condition, and analytic coordinate
covariance. Section~\ref{sec:feedback} states the two repository inputs
and derives the feedback applications. Sections~\ref{sec:growth} and
\ref{sec:holonomic} convert the refined coefficient type into entire
growth and a non-D-finiteness obstruction. Finite verification and the
proposed research program are kept separate in
Sections~\ref{sec:verification} and~\ref{sec:questions}.
\clearpage

\section{Research target, contribution, and dependency boundaries}
\label{sec:context}

\subsection{The precise gap being addressed}

The current ProveIt development separates its transseries library from the
Fabius-function library. The relevant research directory is now
\begin{quote}\small\ttfamily
Analysis/Transseries/docs/series-and-transseries/.
\end{quote}
The repository was inspected at commit \commit. Its transseries README
separates the existing Lean development from the accompanying research
articles \cite{repo,reporeadme}. This article continues the package
\emph{A sharp regularity classification for countable exponential-feedback
transseries} \cite{feedback}, rather than treating another special-function
inverse by a finite coefficient calculation.

That package studies
\begin{equation}
 \U(q)=\sum_{j\ge1}q^j\exp\bigl(\lambda_j\U(q)\bigr),
 \qquad \lambda_j\ge0,\qquad \Q=\U^{-1},
 \label{eq:modelintro}
\end{equation}
where the inverse is compositional. It gives an exact forward Gevrey type
and logarithmic forward coefficient asymptotics, but its README explicitly
limits the reversion theorem to class membership, not the sharp numerical
type. This is a substantive distinction. Knowing that coefficients are
bounded by $CA^n(n!)^s$ for some $A$ does not determine the smallest possible
exponential rate, nor the exact radius of the factorial-normalized Borel
series.

The gap is resolved here by a theorem that is independent of the feedback
model. We preserve the \emph{exact limsup type} under formal inversion for a
large class of coefficient weights, determine precisely when this universal
preservation property holds within the log-convex class, and apply it to
\eqref{eq:modelintro}. The forward coefficients are positive, while inverse
coefficients can have cancellations. The theorem shows that cancellation
cannot lower their weighted exponential limsup.

\subsection{Why a general weight theorem is useful}

A factorial scale alone loses the information relevant to the quadratic
model. Its series have Gevrey order one but zero ordinary Gevrey-one type.
The finer scale
\[
 M_n=\frac{n!}{\log(n+\ee)^{2n}}
\]
has nonzero finite type. Preservation on this scale is substantially more
informative than preservation of Gevrey membership. The same proof applies
to weights growing much more slowly than every positive factorial power,
and also to very rapidly growing weights for which no moderate-growth
property is available.

The proof uses only finite coefficient algebra, log-convexity, and a
subexponential estimate. It does not require an analytic solution of the
feedback equation, a summation operator, positivity of inverse coefficients,
or a description of their signs. This makes the result suitable as an
independent formalization target.

\subsection{What is proved here and what is imported}

All results in Sections~\ref{sec:weights}--\ref{sec:weightspecial} are proved
from elementary formal power-series identities. In
Section~\ref{sec:feedback}, two forward estimates are restated as
\emph{Repository inputs}: the exact forward Gevrey type and the regularly
varying forward coefficient asymptotic. Those two results are used as
mathematical inputs from \cite{feedback};
their full proofs are not independently re-audited in this article.
The inverse consequences and the growth and non-D-finiteness arguments
are proved here from those stated inputs.

The status of each family of assertions is summarized in
Table~\ref{tab:dependencies}. In particular, no application is presented as
Lean-verified merely because it continues a repository that also contains
Lean code.

\begin{table}[htbp]
\centering\small
\begin{tabular}{p{0.37\textwidth}p{0.55\textwidth}}
\toprule
Result & Dependency and status\\
\midrule
Weight concentration; finite inverse majorant & Proved here by finite
inequalities and formal Lagrange inversion.\\[2mm]
Exact type and sharp threshold & Proved here from the finite majorant; no
feedback theorem used.\\[2mm]
Analytic coordinate covariance & Proved here from concentration and the
inverse theorem.\\[2mm]
Forward feedback types and asymptotics & Two explicitly restated results
from the conventional research report \cite{feedback}.\\[2mm]
Inverse feedback types and exact Borel radii & New deductions here from
those inputs and the general inverse theorem.\\[2mm]
Entire growth and non-D-finiteness & Elementary proofs here; application
to feedback uses its imported forward asymptotic.\\[2mm]
Computational evidence & Exact checks at finite orders; neither asymptotic
proof nor formal verification.\\
\bottomrule
\end{tabular}
\caption{Dependency and verification ledger.}
\label{tab:dependencies}
\end{table}

\subsection{Relationship to the literature}

Lagrange inversion is classical; Gessel's survey \cite{gessel} describes
its formal forms and proofs. Stability of ultradifferentiable classes
under composition and inversion is an established subject, including the
weight-sequence and weight-matrix results of Rainer and Schindl
\cite{rainerschindl,stability}. The present object is the exact coefficient
limsup of a \emph{single formal germ}; no realization as an
ultradifferentiable function is presumed. This distinction specifies our
statement, but is not a claim that the literature contains no related or
stronger formulation.

Borel convergence, Borel continuation, and Borel--Laplace summability are
different properties; see Sauzin \cite{sauzin}. We maintain that distinction
throughout. The negative-ray package already in ProveIt \cite{negativeray}
addresses fine negative-direction summation. The present article does not
replace that result with a maximum-modulus assertion or claim angular
summability from a coefficient bound.

The source review was targeted: the current directory inventory, the
transseries README, the classification report's stated theorems and scope,
and the negative-ray report's summary were consulted. It was not a
claim-by-claim audit of the full canonical volume or a complete publication
priority search. The general theorem and its applications are proposed for
independent mathematical review.

\section{Log-convex weights and coefficient type}
\label{sec:weights}

\subsection{Definitions and elementary type rules}

All formal power series are over $\C$, unless otherwise stated. A formal
diffeomorphism is a series $f(z)=a z+O(z^2)$ with $a\ne0$. Such a series
has a unique compositional inverse. A \emph{weight} is a positive sequence
$M=(M_n)_{n\ge0}$ with $M_0=1$.

\begin{definition}
A weight is log-convex if
\begin{equation}
 M_n^2\le M_{n-1}M_{n+1}\qquad(n\ge1).
 \label{eq:logconvex}
\end{equation}
For $n\ge1$, write $r_n=M_n^{1/n}$. The weight is called
superexponential when $r_n\to\infty$.
For $f=\sum_{n\ge0}a_nz^n$, set
\begin{equation}
 \typ(f)=\limsup_{n\to\infty}
       \left(\frac{|a_n|}{M_n}\right)^{1/n}\in[0,\infty].
 \label{eq:type}
\end{equation}
\end{definition}

The choice of $M_n$ refers to \emph{coefficients}, not derivatives. Thus
$M_n=(n!)^s$ gives the factorial-normalized Gevrey type used in this
article. All equalities involving types permit zero and infinity. Positive
finite scalar multiplication of an extended type is interpreted in the
usual way.

\begin{lemma}[Elementary type calculus]
\label{lem:typecalculus}
Finite changes of coefficients or weights do not affect \eqref{eq:type}.
For $c\ne0$,
\[
 \typ(cf)=\typ(f),\qquad \typ(f(cz))=|c|\typ(f).
\]
Furthermore,
\[
 \typ(f+h)\le\max\{\typ(f),\typ(h)\}.
\]
If the two types on the right are unequal, equality holds. If $M$ is
superexponential, every convergent power series has type zero, and adding
any convergent power series preserves the type of an arbitrary formal
series.
\end{lemma}
\begin{proof}
The first statements follow directly by taking $n$th roots. For the sum,
$|a_n+b_n|\le2\max(|a_n|,|b_n|)$ and $2^{1/n}\to1$. If the types are
unequal, write $f=(f+h)-h$ to rule out strict inequality below the larger
one. A convergent series has coefficients bounded by $CR^{-n}$, so its
normalized $n$th root is at most $C^{1/n}R^{-1}/r_n\to0$. When both types
are zero, the sum also has type zero; this completes the last assertion.
\end{proof}

\subsection{The two concentration inequalities}

\begin{lemma}[Log-convex concentration]
\label{lem:concentration}
Let $M$ be log-convex. Then $\mu_n=M_n/M_{n-1}$ and $r_n$ are
nondecreasing. For $0\le\ell\le n$,
\begin{equation}
 \frac{M_{n-\ell}}{M_n}\le r_n^{-\ell}.
 \label{eq:ratio}
\end{equation}
For positive integers $m_1,\ldots,m_k$ with sum $n-1$,
\begin{equation}
 \prod_{i=1}^k M_{m_i+1}
       \le M_{n-k+1}M_2^{k-1}.
 \label{eq:concentrate2}
\end{equation}
For positive integers $j_1,\ldots,j_k$ with sum $n$,
\begin{equation}
 \prod_{i=1}^k M_{j_i}
       \le M_{n-k+1}M_1^{k-1}.
 \label{eq:concentrate1}
\end{equation}
\end{lemma}
\begin{proof}
Log-convexity is equivalent to monotonicity of $\mu_n$. Since
$\log r_n$ is the average of the first $n$ terms of the nondecreasing
sequence $\log\mu_n$, it is nondecreasing. The polygonal function with
values $\log M_n$ at the integers is convex and is zero at the origin.
Its value at $n-\ell$ is therefore at most
$(n-\ell)\log M_n/n$. Exponentiating proves \eqref{eq:ratio}.

For \eqref{eq:concentrate2}, consider the indices $m_i+1\ge2$.
If two indices $a\ge b$ satisfy $b>2$, transfer one unit from $b$ to
$a$. The product cannot decrease, because
\[
 \frac{M_{a+1}M_{b-1}}{M_aM_b}
      =\frac{\mu_{a+1}}{\mu_b}\ge1.
\]
Repeat until all but one index equal two. The remaining index is
$n-k+1$, since the original total is $n-1+k$. This proves
\eqref{eq:concentrate2}. The same argument, with lower bound one in
place of two, proves \eqref{eq:concentrate1}.
\end{proof}

These inequalities do not require $M_1\ge1$ or monotonicity of the
weights themselves. Only positivity and log-convexity are used. The
factor $r_n^{-(k-1)}$ obtained from \eqref{eq:ratio} is the mechanism that
suppresses the combinatorial multiplicity of nonlinear inverse terms.

\section{A finite coefficient bound for formal reversion}
\label{sec:majorant}

\subsection{The coefficient formula}

Write
\begin{equation}
 f(z)=z+\sum_{j\ge2}a_jz^j=z(1+h(z)),\qquad
 g(z)=f^{-1}(z)=z+\sum_{n\ge2}b_nz^n.
 \label{eq:fg}
\end{equation}
Lagrange inversion gives, for $n\ge1$,
\begin{equation}
 b_n=\frac1n[z^{n-1}](1+h(z))^{-n}.
 \label{eq:lagrange}
\end{equation}
A formal residue proof is included in Appendix~\ref{app:lagrange}; no
analytic convergence is required. Expanding the negative integer power
in \eqref{eq:lagrange}, for $n\ge2$ one obtains the finite identity
\begin{equation}
 b_n=\frac1n\sum_{k=1}^{n-1}(-1)^k\binom{n+k-1}{k}
 \sum_{\substack{m_1+\cdots+m_k=n-1\\m_i\ge1}}
       \prod_{i=1}^k a_{m_i+1}.
 \label{eq:inversefinite}
\end{equation}
The term $k=1$ equals $-a_n$. The other terms can be large and can cancel;
it would be invalid to infer a lower bound for $|b_n|$ by keeping only
that first term. Our proof instead obtains a one-sided upper bound that
can be applied twice, once in each direction of inversion.

\begin{theorem}[Finite inverse majorant]
\label{thm:finite}
Assume $M$ is log-convex and
\begin{equation}
 |a_j|\le CA^jM_j\qquad(j\ge2),\qquad C,A>0.
 \label{eq:inputbound}
\end{equation}
Set $D=CA$. For every $n\ge2$,
\begin{equation}
 |b_n|\le\frac{A^{n-1}}n
 \sum_{k=1}^{n-1}\binom{n+k-1}{k}\binom{n-2}{k-1}
       D^kM_{n-k+1}M_2^{k-1}.
 \label{eq:finitemajorant}
\end{equation}
In particular, with
\begin{equation}
 X_n=\frac{2DM_2n^2}{r_n},
 \label{eq:X}
\end{equation}
one has the explicit bound
\begin{align}
 |b_n|&\le2DA^{n-1}M_n
       \sum_{\ell=0}^{n-2}\frac{X_n^\ell}{\ell!(\ell+1)!}
 \notag\\
 &\le2DA^{n-1}M_n\exp\left(2n\sqrt{\frac{2DM_2}{r_n}}\right).
 \label{eq:explicitmajorant}
\end{align}
\end{theorem}
\begin{proof}
The number of positive ordered compositions of $n-1$ into $k$ parts
is $\binom{n-2}{k-1}$. Each summand in the inner sum of
\eqref{eq:inversefinite} is bounded by
\[
 C^k A^{n-1+k}\prod_iM_{m_i+1}
 \le A^{n-1}D^kM_{n-k+1}M_2^{k-1},
\]
by Lemma~\ref{lem:concentration}. Taking absolute values proves
\eqref{eq:finitemajorant}.

For $1\le k\le n-1$, use
\[
 \binom{n+k-1}{k}\le\frac{(2n)^k}{k!},\qquad
 \binom{n-2}{k-1}\le\frac{n^{k-1}}{(k-1)!},\qquad
 \frac{M_{n-k+1}}{M_n}\le r_n^{-(k-1)}.
\]
Substitution in \eqref{eq:finitemajorant}, followed by $\ell=k-1$,
gives the first inequality in \eqref{eq:explicitmajorant}.
For $X\ge0$,
\begin{align*}
 \sum_{\ell\ge0}\frac{X^\ell}{\ell!(\ell+1)!}
 &\le\sum_{\ell\ge0}\frac{X^\ell}{(\ell!)^2}\\
 &=\sum_{\ell\ge0}\binom{2\ell}{\ell}\frac{X^\ell}{(2\ell)!}
 \le\sum_{\ell\ge0}\frac{(2\sqrt X)^{2\ell}}{(2\ell)!}
 \le\exp(2\sqrt X).
\end{align*}
This proves the second inequality.
\end{proof}

\subsection{An effective zero-loss interpretation}

The finite sum \eqref{eq:finitemajorant} is an exact-arithmetic bound when
the data are rational. The exponential simplification is more convenient
for asymptotics. In particular,
\begin{equation}
 \left(\frac{|b_n|}{M_n}\right)^{1/n}
 \le A\left(\frac{2D}{A}\right)^{1/n}
       \exp\left(2\sqrt{\frac{2DM_2}{r_n}}\right).
 \label{eq:nthrootbound}
\end{equation}
Every factor beyond $A$ tends to one if $r_n\to\infty$.

For a concrete tolerance $\varepsilon>0$, the condition
\[
 r_n\ge\frac{8DM_2}{\log(1+\varepsilon)^2}
\]
makes the exponential factor in \eqref{eq:nthrootbound} at most
$1+\varepsilon$. A separate finite lower bound on $n$ makes
$(2D/A)^{1/n}\le1+\varepsilon$ as well. Thus the theorem gives a
computable onset for an inverse rate $A(1+\varepsilon)^2$, provided a
verified input bound and growth bound for $r_n$ are supplied. A mere
limsup statement need not supply those bounds effectively.

\section{Exact type invariance and its sharp boundary}
\label{sec:invariance}

\begin{theorem}[Zero-loss reversion]
\label{thm:zeroloss}
Let $M$ be positive and log-convex, with $M_0=1$ and $r_n\to\infty$.
For every formal diffeomorphism $f(z)=az+O(z^2)$,
\begin{equation}
 \boxed{\displaystyle
 \typ(f^{-1})=\frac{\typ(f)}{|a|}.}
 \label{eq:zeroloss}
\end{equation}
In particular, tangent-to-identity inversion preserves type exactly,
including zero and infinite type.
\end{theorem}
\begin{proof}
First suppose $a=1$ and $\typ(f)<\infty$. For every $A>\typ(f)$,
definition \eqref{eq:type} gives a constant $C>0$ such that
\eqref{eq:inputbound} holds at every $j\ge2$: finitely many initial
coefficients are absorbed in $C$. Theorem~\ref{thm:finite} and
\eqref{eq:nthrootbound} give $\typ(f^{-1})\le A$. Letting $A$ decrease
to $\typ(f)$ proves
\[
 \typ(f^{-1})\le\typ(f).
\]
This includes $\typ(f)=0$, since every $A>0$ may be used. The inverse
also has finite type, so apply the same argument to $f^{-1}$, whose
inverse is $f$, to obtain the reverse inequality.

If $\typ(f)=\infty$ but $\typ(f^{-1})<\infty$, the already proved
finite-type upper bound applied to $f^{-1}$ would imply
$\typ(f)<\infty$, a contradiction. This handles the infinite case.

For arbitrary $a\ne0$, set $F=f/a$. Output scaling does not change
type and $F'(0)=1$. Since
\[
 F^{-1}(z)=f^{-1}(az),
\]
Lemma~\ref{lem:typecalculus} and the tangent case give
$|a|\typ(f^{-1})=\typ(F^{-1})=\typ(F)=\typ(f)$.
\end{proof}

\begin{theorem}[Necessary and sufficient weight condition]
\label{thm:iff}
For a positive log-convex weight $M_0=1$, the following are equivalent:
\begin{enumerate}[label=(\roman*)]
\item $M_n^{1/n}\to\infty$.
\item Every tangent-to-identity formal diffeomorphism satisfies
$\typ(f^{-1})=\typ(f)$.
\item The inverse of every tangent-to-identity polynomial has type zero.
\end{enumerate}
\end{theorem}
\begin{proof}
The implication (i)$\Rightarrow$(ii) is Theorem~\ref{thm:zeroloss},
and (ii)$\Rightarrow$(iii) follows because polynomials have type zero.
For the converse, the roots $r_n$ are positive and nondecreasing, so
failure of (i) means $r_n\to L$ for some $0<L<\infty$. Take
$f(z)=z-z^2$. Its inverse is
\begin{equation}
 f^{-1}(z)=\frac{1-\sqrt{1-4z}}2
 =\sum_{n\ge1}\frac1n\binom{2n-2}{n-1}z^n.
 \label{eq:catalan}
\end{equation}
The coefficients have $n$th-root limit $4$, by Stirling's formula.
Consequently $\typ(f^{-1})=4/L>0$, contradicting (iii).
\end{proof}

\begin{remark}[The analytic endpoint is genuinely different]
For $M_n=1$, type is the reciprocal of the ordinary radius of
convergence. In \eqref{eq:catalan}, a polynomial of type zero has an
inverse of type four. Thus no version of exact type preservation that
includes all analytic radii can hold. Superexponential normalization
makes every analytic germ type zero and removes this analytic
singularity obstruction from the coefficient scale.
\end{remark}

\subsection{What exact type does and does not say}

If $0<\typ(f)=T<\infty$, the equality in
Theorem~\ref{thm:zeroloss} says that no uniform exponential improvement
of the inverse coefficients is possible. For every $0<\varepsilon<T$,
infinitely many indices satisfy
\[
 |b_n|\ge (T-\varepsilon)^nM_n
\]
in the tangent case, while for every $\varepsilon>0$ all sufficiently
large indices satisfy
\[
 |b_n|\le (T+\varepsilon)^nM_n.
\]
These statements are compatible with zeros, sign changes, sparse large
coefficients, and substantial subexponential cancellation. They do not
imply that $|b_n|^{1/n}/r_n$ has a limit. Reversion twice is what proves
the sharp lower limsup; no unproved sign assumption enters.

\section{Exact covariance under analytic coordinate changes}
\label{sec:coordinates}

\subsection{Postcomposition: a second finite estimate}

\begin{proposition}[Analytic outer functions]
\label{prop:post}
Assume $M$ satisfies the hypotheses of Theorem~\ref{thm:zeroloss}.
Let $f(0)=0$ be any formal series and let $H$ be analytic near zero.
Then
\begin{equation}
 \typ(H\circ f)\le\typ(f).
 \label{eq:post}
\end{equation}
If $H'(0)\ne0$, equality holds.
More explicitly, suppose
\[
 |[z^j]f|\le CA^jM_j\quad(j\ge1),\qquad
 |[w^k]H|\le BR^{-k}\quad(k\ge1).
\]
For $n\ge1$,
\begin{equation}
 |[z^n]H(f(z))|
 \le\frac{BC}{R}A^nM_n
 \exp\left(\frac{CM_1n}{Rr_n}\right).
 \label{eq:postbound}
\end{equation}
\end{proposition}
\begin{proof}
There are $\binom{n-1}{k-1}$ positive compositions of $n$ into $k$
parts. Lemma~\ref{lem:concentration} gives
\begin{align*}
 |[z^n]H(f(z))|
 &\le B A^n\sum_{k=1}^n (C/R)^k
    \binom{n-1}{k-1}M_{n-k+1}M_1^{k-1}\\
 &\le\frac{BC}{R}A^nM_n
    \left(1+\frac{CM_1}{Rr_n}\right)^{n-1},
\end{align*}
which implies \eqref{eq:postbound}. For finite $\typ(f)$, choose any
$A>\typ(f)$ and absorb all initial coefficients in $C$. Taking $n$th
roots and then the limsup proves \eqref{eq:post}; if the right side is
infinite there is nothing to prove. If $H'(0)\ne0$, the analytic inverse
germ of $H-H(0)$ exists. Apply the upper inequality to this inverse and
to $H(f)-H(0)$, then use invariance under finite coefficient changes.
\end{proof}

\subsection{Precomposition and conjugacy}

\begin{theorem}[Analytic coordinate covariance]
\label{thm:coordinates}
Let $M$ satisfy the hypotheses of Theorem~\ref{thm:zeroloss}.
Let $f$ be any formal series, let $g$ be analytic with $g(0)=0$ and
$g'(0)=c\ne0$, and let $H$ be analytic near $f(0)$ with
$H'(f(0))\ne0$. Then
\begin{equation}
 \boxed{\displaystyle
 \typ(H\circ f\circ g)=|c|\typ(f).}
 \label{eq:coordinates}
\end{equation}
In particular, conjugacy by an analytic tangent-to-identity change of
variable preserves exact type.
\end{theorem}
\begin{proof}
Start with a formal diffeomorphism $f$ and write $a=f'(0)\ne0$.
The inverse factorization is
\[
 (f\circ g)^{-1}=g^{-1}\circ f^{-1}.
\]
The germ $g^{-1}$ is analytic and has nonzero derivative, so
Proposition~\ref{prop:post} and Theorem~\ref{thm:zeroloss} give
\[
 \frac{\typ(f\circ g)}{|ac|}
 =\typ((f\circ g)^{-1})
 =\typ(f^{-1})
 =\frac{\typ(f)}{|a|}.
\]
Thus $\typ(f\circ g)=|c|\typ(f)$.

For general $f$, subtract its constant term and choose a linear
polynomial $\ell$ so that $F=f-f(0)+\ell$ has nonzero linear
coefficient. Both $F-f$ and $F\circ g-f\circ g$ are analytic.
Lemma~\ref{lem:typecalculus} therefore transfers the equality from $F$
to $f$. Finally apply Proposition~\ref{prop:post} to the analytic outer
function translated to the constant term $f(0)$.
\end{proof}

The theorem does not assert a type formula for composition of two
arbitrary divergent series. Analyticity of the coordinate maps is an
essential part of the statement proved here. It also does not assert
invariance under changes with zero derivative, such as ramification
$z\mapsto z^d$; those changes rescale coefficient indices rather than
merely the exponential type.

\section{Gevrey and logarithmically refined weights}
\label{sec:weightspecial}

\subsection{A family extending all positive Gevrey orders}

For real parameters $s,\beta$, define for $n\ge1$
\begin{equation}
 W_n(s,\beta)=(n!)^s\bigl(\log(n+\ee)\bigr)^{\beta n},
 \qquad W_0=1.
 \label{eq:logweight}
\end{equation}
The logarithmic shift ensures positivity at small indices; changing it
does not affect any type considered below.

\begin{lemma}[Admissibility of the refined weights]
\label{lem:refinedweight}
Suppose either $s>0$ and $\beta\in\R$, or $s=0$ and $\beta>0$.
After changing only finitely many initial terms, \eqref{eq:logweight}
is a log-convex superexponential weight. Moreover,
\begin{equation}
 W_n(s,\beta)^{1/n}\sim(n/\ee)^s(\log n)^\beta.
 \label{eq:weightroot}
\end{equation}
\end{lemma}
\begin{proof}
Write $F_n=\log W_n$ and $B(x)=x\log\log(x+\ee)$. Then
\[
 F_{n+1}-2F_n+F_{n-1}
 =s\log\frac{n+1}{n}+\beta\bigl(B(n+1)-2B(n)+B(n-1)\bigr).
\]
Elementary differentiation gives
\[
 B''(x)=\frac1{x\log x}\left(1+O((\log x)^{-1})\right).
\]
The discrete second difference of $B$ is the integral of $B''$ against
a nonnegative triangular kernel of total mass one on $[n-1,n+1]$.
Consequently the second difference of $F_n$ is
\[
 \frac{s}{n}+O(n^{-2})+
 \frac{\beta}{n\log n}\left(1+O((\log n)^{-1})\right),
\]
which is eventually positive under either hypothesis. Moreover, using
Stirling's formula only for $\log(n!)$, for $s>0$ one gets
\[
 F_{n+1}-F_n-\frac{F_n}{n}=s+o(1)>0
\]
for sufficiently large $n$. For $s=0$, $\beta>0$, the eventual
convexity of $B$ and the exact identity
\[
 B'(x)-\frac{B(x)}x=\frac{x}{(x+\ee)\log(x+\ee)}>0
\]
give the same positivity directly. Thus there is an integer $N$ in
the convex tail such that
\[
 \frac{F_N}{N}\le F_{N+1}-F_N.
\]
Keep $W_n$ for $n\ge N$, and set $\widetilde M_n=W_N^{n/N}$ for
$0\le n<N$. The initial logarithms lie on a straight line; the displayed
inequality makes their join with the convex tail convex. Thus
$\widetilde M$ is a globally log-convex weight with the same tail.
Stirling's formula proves \eqref{eq:weightroot}, which tends to infinity
under the stated hypotheses.
\end{proof}

\begin{corollary}[Exact refined-type inversion]
\label{cor:refinedgeneral}
Under either hypothesis in Lemma~\ref{lem:refinedweight}, let
\[
 T_{s,\beta}(f)=\limsup_n
 \left(\frac{|[z^n]f|}{(n!)^s\log(n+\ee)^{\beta n}}\right)^{1/n}.
\]
For every formal diffeomorphism with derivative $a\ne0$,
\begin{equation}
 T_{s,\beta}(f^{-1})=\frac{T_{s,\beta}(f)}{|a|}.
 \label{eq:refinedgeneral}
\end{equation}
The analytic coordinate law \eqref{eq:coordinates} holds for these
refined types as well.
\end{corollary}
\begin{proof}
Use the finite modification from Lemma~\ref{lem:refinedweight} in
Theorems~\ref{thm:zeroloss} and \ref{thm:coordinates}. Finite changes
of weights do not alter type.
\end{proof}

For $s>0$, $\beta=0$, this gives exact Gevrey-type preservation, not
merely closure of the Gevrey class. For $s=0$, $\beta>0$, it gives the
same exact preservation on a scale slower than every positive Gevrey
weight. By contrast, $s=\beta=0$ is the analytic endpoint excluded by
Theorem~\ref{thm:iff}.

\subsection{Size of the explicit overhead}

Fix an admissible globally log-convex modification and a bound
\eqref{eq:inputbound}. From \eqref{eq:explicitmajorant} and
\eqref{eq:weightroot},
\begin{equation}
 \log\frac{|b_n|}{A^{n-1}M_n}
 \le O(1)+O\left(n^{1-s/2}(\log n)^{-\beta/2}\right),
 \label{eq:overheadrefined}
\end{equation}
with the convention that a zero coefficient satisfies the inequality.
The constants depend on the fixed input bound and the finite weight
modification. In the sub-Gevrey case $s=0$, $\beta>0$, the overhead is
$O(n/(\log n)^{\beta/2})$, still $o(n)$.

This is a \emph{proved upper bound}, not a sharp asymptotic for nonlinear
inverse corrections. Its role is to retain the exponential constant, and
its optimization is a separate research question.

\section{Application to countable exponential feedback}
\label{sec:feedback}

\subsection{Formal existence and a positive coefficient formula}

Fix finite real numbers $\lambda_j\ge0$, without any growth assumption.
Consider
\begin{equation}
 \Phi(q,u)=\sum_{j\ge1}q^j e^{\lambda_j u},\qquad
 \U(q)=\Phi(q,\U(q)).
 \label{eq:feedback}
\end{equation}
The coefficient of $q^n$ on the right depends only on coefficients of
$\U$ below degree $n$. Therefore the equation determines a unique series
$\U\in q\R[[q]]$, with $u_1=1$ if
$\U=\sum_{n\ge1}u_nq^n$. The countable kernel need not converge in a
complex neighborhood of $(0,0)$; formal existence does not require such
convergence.

\begin{proposition}[Finite positive Lagrange formula]
\label{prop:feedbackcoef}
Let $\mathbf m=(m_1,m_2,\ldots)$ have finite support, and write
$|\mathbf m|=\sum m_j$, $w(\mathbf m)=\sum jm_j$. Then
\begin{equation}
 u_n=\sum_{w(\mathbf m)=n}
 \frac{\bigl(\sum_j\lambda_jm_j\bigr)^{|\mathbf m|-1}}
      {\prod_jm_j!},\qquad n\ge1,
 \label{eq:positiveformula}
\end{equation}
where $0^0=1$. In particular $u_n\ge1$.
\end{proposition}
\begin{proof}
Introduce a marker $t$ and solve $V=t\Phi(q,V)$ formally. Lagrange
inversion in $t$ gives
\[
 [t^k]V=\frac1k[u^{k-1}]\Phi(q,u)^k.
\]
Every term has $q$-order at least $k$, so setting $t=1$ is coefficientwise
well-defined. For an ordered $k$-tuple $(j_1,\ldots,j_k)$, the contribution
is
\[
 \frac1{k!}\bigl(\lambda_{j_1}+\cdots+\lambda_{j_k}\bigr)^{k-1}
 q^{j_1+\cdots+j_k}.
\]
Grouping tuples by their multiplicities gives
\eqref{eq:positiveformula}. All terms are nonnegative, and the tuple of
length one with entry $n$ contributes one.
\end{proof}

With $q=e^{-x}$, the formal series becomes
$\sum_{n\ge1}u_ne^{-nx}$. It already has a simple well-based
one-generator exponential support. The questions below concern the
coefficient growth on that support, not a failure of support
well-ordering. Since $u_1=1$, the inverse $\Q(u)=\sum_{n\ge1}q_nu^n$
exists formally and satisfies the literal formal identity
\begin{equation}
 u=\sum_{j\ge1}\Q(u)^j e^{\lambda_j u}.
 \label{eq:inversekernel}
\end{equation}

\subsection{The two forward inputs}

The following are the precise results taken from \cite{feedback}. They
are labeled as inputs to avoid attributing their proofs to the present
article.

\begin{inputtheorem}[Exact forward Gevrey type]
\label{input:forwardtype}
For the nonnegative-slope model \eqref{eq:feedback} and every $s>0$,
\begin{equation}
 T_s(\U):=\limsup_n\left(\frac{u_n}{(n!)^s}\right)^{1/n}
 =\left(\frac{s+1}{s}\right)^{s+1}
   \limsup_{j\to\infty}\frac{\lambda_j}{(j\log j)^{s+1}}.
 \label{eq:forwardtype}
\end{equation}
The equality holds in $[0,\infty]$.
\end{inputtheorem}

\begin{inputtheorem}[Regularly varying forward slopes]
\label{input:forwardasymptotic}
Suppose $a>0$, $p>1$, $\beta\in\R$, and
$\lambda_j\sim a j^p(\log j)^\beta$. Then
\begin{align}
 \log u_n={}&(p-1)n\log n+(\beta-p)n\log\log n
 \notag\\
 &+\left(\log a+p\log\frac{p}{p-1}-(p-1)\right)n+o(n).
 \label{eq:forwardasymptotic}
\end{align}
\end{inputtheorem}

In the cited source these are the statements labeled
\texttt{thm:classification}, equation \texttt{eq:exacttype}, and
\texttt{thm:regular}, equation \texttt{eq:logasymptotic}, respectively.
They are conventional research results, not formalized statements.
Neither our abstract inverse theorem nor its coordinate-change theorem
uses these inputs; the remaining feedback applications do.

\subsection{The sharp inverse classification}

\begin{theorem}[Exact inverse Gevrey type]
\label{thm:inversefeedback}
For \eqref{eq:feedback} and its formal inverse, for every $s>0$,
\begin{equation}
 \boxed{\displaystyle
 T_s(\Q)=T_s(\U)
 =\left(\frac{s+1}{s}\right)^{s+1}
   \limsup_{j\to\infty}\frac{\lambda_j}{(j\log j)^{s+1}}.}
 \label{eq:inversefeedback}
\end{equation}
Consequently $\Q$ belongs to Gevrey class $s$ exactly when
$\lambda_j=O((j\log(j+1))^{s+1})$, and its exact type is the constant
in \eqref{eq:inversefeedback}.
\end{theorem}
\begin{proof}
For $M_n=(n!)^s$, the weight is log-convex and superexponential.
Theorem~\ref{thm:zeroloss} gives $T_s(\Q)=T_s(\U)$ because $u_1=1$.
Insert Repository input~\ref{input:forwardtype}. A finite type is
equivalent to a bound $CA^n(n!)^s$, and the limsup on the right is
finite exactly under the stated big-$O$ condition, since every individual
slope is finite.
\end{proof}

The class-membership equivalence was already available in
\cite{feedback}; the addition in \eqref{eq:inversefeedback} is the exact
constant for the inverse, including the zero and infinite cases.
Nonnegativity is an assumption on the primitive slopes and the imported
forward results, not on the inverse coefficients.

\begin{corollary}[Exact inverse Borel radius]
\label{cor:borelradius}
Define the integrated ordinary Borel series by
\begin{equation}
 (\B\Q)(\zeta)=\sum_{n\ge1}\frac{q_n}{n!}\zeta^n.
 \label{eq:boreldef}
\end{equation}
Writing
\[
 A_*:=\limsup_{j\to\infty}\frac{\lambda_j}{(j\log j)^2},
\]
its exact radius of convergence is
\begin{equation}
 \boxed{\displaystyle R_{\B\Q}=\frac1{4A_*},}
 \label{eq:borelradius}
\end{equation}
with $1/0=\infty$ and $1/\infty=0$.
More generally, the series
$\sum q_n\zeta^n/(n!)^s$ has radius $1/T_s(\Q)$.
\end{corollary}
\begin{proof}
Apply Cauchy--Hadamard to \eqref{eq:boreldef}, followed by
Theorem~\ref{thm:inversefeedback} with $s=1$. The general assertion
is the same root test with the stated normalization.
\end{proof}

For example, if $\lambda_j\sim A(j\log j)^2$ with $A>0$, the exact
radius is $1/(4A)$. If $\lambda_j=j^2$, the Borel transform is entire.
If $A_*=\infty$, it has zero radius as a formal Borel series.
A finite positive radius implies at least one obstruction to analyticity
on the boundary circle, but does not identify its angle or its local
nature. An infinite radius says nothing by itself about a Laplace
integral on any prescribed ray. For nonintegral $s$, the factorial
normalization above is the one explicitly defined here, and should not
be confused with a different $k$-summability convention involving
$\Gamma(1+sn)$.

\subsection{A sharp logarithmically refined inverse invariant}

\begin{theorem}[Regularly varying inverse type]
\label{thm:refinedfeedback}
Under Repository input~\ref{input:forwardasymptotic}, put
\begin{equation}
 s=p-1,\qquad \gamma=\beta-p,\qquad
 K=a\left(\frac{p}{p-1}\right)^p.
 \label{eq:parameters}
\end{equation}
Then
\begin{equation}
 \boxed{\displaystyle
 \limsup_{n\to\infty}
 \left(\frac{|q_n|}{(n!)^{p-1}\log(n+\ee)^{(\beta-p)n}}\right)^{1/n}
 =K.}
 \label{eq:refinedfeedback}
\end{equation}
For the forward coefficients, the corresponding normalized $n$th roots
have an actual limit equal to $K$.
\end{theorem}
\begin{proof}
Stirling's formula gives
\[
 \log\bigl((n!)^{p-1}\log(n+\ee)^{(\beta-p)n}\bigr)
 =(p-1)n\log n-(p-1)n
   +(\beta-p)n\log\log n+o(n).
\]
Subtract this from \eqref{eq:forwardasymptotic} and divide by $n$.
The resulting limit is $\log K$. Thus the forward refined type is $K$,
and in fact the forward normalized roots converge to it. The weight
satisfies Lemma~\ref{lem:refinedweight} because $p-1>0$.
Corollary~\ref{cor:refinedgeneral} transfers its exact type to $\Q$.
\end{proof}

An immediate consequence is a two-sided exponential envelope in the
limsup sense: for each $\varepsilon>0$,
\[
 |q_n|\le(K+\varepsilon)^n(n!)^{p-1}
             \log(n+\ee)^{(\beta-p)n}
\]
for all sufficiently large $n$; if $0<\varepsilon<K$, the reversed bound
with $K-\varepsilon$ holds at infinitely many $n$. Modifying finitely
many slopes does not change this invariant when the same regularly
varying tail is retained. The theorem does not say that such a
modification preserves each coefficient or each analytic summation.

\begin{corollary}[Quadratic feedback]
\label{cor:quadratic}
For $\lambda_j=j^2$,
\begin{equation}
 \limsup_{n\to\infty}
 \left(\frac{|q_n|\log(n+\ee)^{2n}}{n!}\right)^{1/n}=4.
 \label{eq:quadratic}
\end{equation}
Both $\U$ and $\Q$ have Gevrey-one type zero, belong to no Gevrey
class $s<1$ with $s>0$, and have entire ordinary Borel transforms.
\end{corollary}
\begin{proof}
Use $a=1,p=2,\beta=0$ in Theorem~\ref{thm:refinedfeedback}.
The exact ordinary Gevrey types also follow from
\eqref{eq:inversefeedback}: the slope ratio tends to zero for $s=1$
and to infinity for $0<s<1$. The Borel conclusion follows from
Corollary~\ref{cor:borelradius}.
\end{proof}

\section{Entire Borel growth from the refined type}
\label{sec:growth}

\subsection{A coefficient-to-growth lemma that permits cancellation}

For an entire function $F$, write
$\mathcal M_F(r)=\max_{|z|=r}|F(z)|$.
In iterated logarithms, replacing $\mathcal M_F(r)$ by
$\max(\ee,\mathcal M_F(r))$ at small radii makes all expressions
defined and does not affect any limit below.

\begin{lemma}[Iterated-logarithmic growth]
\label{lem:entiregrowth}
Suppose $F(z)=\sum_{n\ge0}c_nz^n$, $d>0$, $K>0$, and
\begin{equation}
 \limsup_{n\to\infty}|c_n|^{1/n}\log(n+\ee)^d=K.
 \label{eq:rootdecay}
\end{equation}
Then $F$ is entire and
\begin{equation}
 \limsup_{r\to\infty}
 \frac{\log\log\mathcal M_F(r)}{r^{1/d}}=K^{1/d}.
 \label{eq:entiregrowth}
\end{equation}
In particular $F$ has infinite classical entire order.
\end{lemma}
\begin{proof}
Equation~\eqref{eq:rootdecay} implies $\limsup|c_n|^{1/n}=0$,
so $F$ is entire. For an upper bound choose $L>K$ and $\eta>0$.
For all sufficiently large $n$,
\[
 |c_n|\le\frac{L^n}{\log(n+\ee)^{dn}}.
\]
For sufficiently large $r$, set
\[
 N=\left\lceil\exp\bigl((1+\eta)(Lr)^{1/d}\bigr)\right\rceil.
\]
For $n\ge N$, one has
$Lr/\log(n+\ee)^d\le(1+\eta)^{-d}=:\rho<1$.
The tail is at most $\sum_{n\ge N}\rho^n$ in absolute value on
$|z|=r$. The terms before $N$ are at most a fixed polynomial in $r$
plus $N\max(1,Lr)^N$. Consequently
\[
 \log\log\mathcal M_F(r)
 \le (1+\eta)(Lr)^{1/d}+O(\log\log r).
\]
Taking the limsup and then $L\downarrow K$, $\eta\downarrow0$
gives the upper bound in \eqref{eq:entiregrowth}.

For the lower bound, choose $0<\varepsilon<K$ and $\eta>0$.
There are infinitely many $n$ for which
\[
 |c_n|\ge\frac{(K-\varepsilon)^n}{\log(n+\ee)^{dn}}.
\]
At the radii
\[
 r_n=\frac{\ee^{d\eta}\log(n+\ee)^d}{K-\varepsilon},
\]
Cauchy's coefficient inequality gives
$\mathcal M_F(r_n)\ge|c_n|r_n^n\ge\ee^{d\eta n}$.
It follows that
\[
 \limsup_{r\to\infty}
 \frac{\log\log\mathcal M_F(r)}{r^{1/d}}
 \ge \ee^{-\eta}(K-\varepsilon)^{1/d}.
\]
Let $\eta\downarrow0$ and $\varepsilon\downarrow0$.
The positive finite value in \eqref{eq:entiregrowth} implies
$\limsup\log\log\mathcal M_F(r)/\log r=\infty$, which is infinite
classical order.
\end{proof}

No positivity assumption is present in this lemma. The lower bound uses
Cauchy's inequality for a coefficient, rather than a value of the
function on a prescribed ray. This distinction is particularly important
for a signed inverse series.

\subsection{Feedback transforms of integer factorial order}

\begin{theorem}[Exact inverse Borel maximum-modulus type]
\label{thm:borelgrowth}
Assume
\[
 \lambda_j\sim a j^p(\log j)^\beta,\qquad
 p=m+1,\quad m\in\N,\ m\ge1,\quad a>0,\quad\beta<p.
\]
Set $d=p-\beta>0$ and $K=a(p/(p-1))^p$. Define
\begin{equation}
 B_{m,Q}(z)=\sum_{n\ge1}\frac{q_n}{(n!)^m}z^n,\qquad
 B_{m,U}(z)=\sum_{n\ge1}\frac{u_n}{(n!)^m}z^n.
 \label{eq:Bmq}
\end{equation}
Both are entire, and each satisfies
\begin{equation}
 \limsup_{r\to\infty}
 \frac{\log\log\mathcal M_{B}(r)}{r^{1/d}}
 =K^{1/d}.
 \label{eq:Bmgrowth}
\end{equation}
\end{theorem}
\begin{proof}
Theorem~\ref{thm:refinedfeedback}, with $m=p-1$, says that the
coefficients of $B_{m,Q}$ satisfy \eqref{eq:rootdecay}. For $B_{m,U}$
the normalized root even has a limit, so the same hypothesis holds.
Apply Lemma~\ref{lem:entiregrowth} to each transform.
\end{proof}

For $m=1$ this is the ordinary integrated Borel transform; for larger
integer $m$ it is the explicitly defined repeated factorial
normalization, not an unstated alternative convention. In the quadratic
case, \eqref{eq:Bmgrowth} becomes
\begin{equation}
 \boxed{\displaystyle
 \limsup_{r\to\infty}
 \frac{\log\log\mathcal M_{\B\Q}(r)}{\sqrt r}=2.}
 \label{eq:quadraticgrowth}
\end{equation}

The coefficient information proved here justifies a \emph{limsup} for
the inverse. It does not justify replacing that limsup by a limit or
$\mathcal M_{\B\Q}(r)$ by $|(\B\Q)(r)|$. In particular, infinite
maximum-modulus order does not contradict summability on a direction
where the same entire function grows much more slowly. This is
consistent with the separate negative-direction result in
\cite{negativeray}.

\section{Non-D-finiteness and analytic-normal-form obstructions}
\label{sec:holonomic}

\subsection{Two elementary holonomic facts}

\begin{definition}
A formal series $F(z)$ is D-finite if it satisfies a nonzero linear
differential equation
\begin{equation}
 p_r(z)F^{(r)}(z)+\cdots+p_0(z)F(z)=0,
 \qquad p_j\in\C[z],\quad p_r\ne0.
 \label{eq:Dfinite}
\end{equation}
A sequence is P-recursive if, for all sufficiently large $n$, it
satisfies a nonzero linear recurrence with polynomial coefficients in $n$.
\end{definition}

\begin{lemma}[Factorial division preserves D-finiteness]
\label{lem:factorialD}
If $F(z)=\sum a_nz^n$ is a D-finite formal series and $m$ is a
nonnegative integer, then $\sum a_nz^n/(n!)^m$ is D-finite formally.
\end{lemma}
\begin{proof}
Coefficient extraction in \eqref{eq:Dfinite} gives a polynomial-coefficient
recurrence for $a_n$, after shifting its indices to be nonnegative.
Write it as
\[
 \sum_{j=0}^rP_j(n)a_{n+j}=0
\]
for all sufficiently large $n$, with not all $P_j$ zero. Put
$a_n=(n!)^m c_n$ and divide by $(n!)^m$. This gives
\[
 \sum_{j=0}^r P_j(n)\bigl((n+1)\cdots(n+j)\bigr)^m c_{n+j}=0.
\]
Every coefficient is polynomial because $m$ is an integer, and the
recurrence is nonzero. Such a recurrence yields a linear differential
equation for $\sum c_nz^n$ by multiplying by powers of $z$, summing in
$n$, and replacing $n$ by the Euler operator $z\,\dd/\dd z$.
Finitely many exceptional initial indices contribute a polynomial
inhomogeneous term, which is annihilated by a sufficiently high
derivative. This yields a nonzero homogeneous linear differential
operator, proving D-finiteness.
\end{proof}

\begin{lemma}[Entire D-finite functions have finite order]
\label{lem:finiteorder}
Every entire D-finite function has finite classical entire order.
\end{lemma}
\begin{proof}
The zero function is harmless. A nonzero entire solution has an equation
\eqref{eq:Dfinite} of order $r\ge1$. Choose $R$ so large that all zeros
of $p_r$ lie in $|z|<R$. Outside that disk, the vector
\[
 Y=(F,F',\ldots,F^{(r-1)})^{\mathsf T}
\]
satisfies a companion system $Y'=A(z)Y$ whose entries are rational
functions. There exist $C>0$ and an integer $L\ge0$ with
$\|A(z)\|\le C|z|^L$ for $|z|\ge R$. The initial values
$Y(R\ee^{i\theta})$ are uniformly bounded in $\theta$, since all
components are entire. Along a radial segment, Gronwall's inequality
therefore gives, uniformly in $\theta$,
\[
 \|Y(r\ee^{i\theta})\|
 \le C_0\exp\left(C\int_R^r t^L\,\dd t\right)
 \le C_0\exp(C_1r^{L+1}).
\]
Thus $\mathcal M_F(r)\le C_0\exp(C_1r^{L+1})$, which is finite
classical order.
\end{proof}

\subsection{An obstruction valid for signed coefficient sequences}

\begin{theorem}[Logarithmic factorial obstruction]
\label{thm:obstruction}
Let $m\ge1$ be an integer, $d>0$, and $K>0$. Suppose a formal series
$F(z)=\sum a_nz^n$ satisfies
\begin{equation}
 \limsup_{n\to\infty}
 \left(\frac{|a_n|\log(n+\ee)^{dn}}{(n!)^m}\right)^{1/n}=K.
 \label{eq:obstructionhyp}
\end{equation}
Then $F$ is not D-finite, and its coefficient sequence is not
P-recursive.
\end{theorem}
\begin{proof}
Lemma~\ref{lem:entiregrowth} says that
$B(z)=\sum a_nz^n/(n!)^m$ is entire and of infinite order. If $F$
were D-finite, Lemma~\ref{lem:factorialD} would make $B$ D-finite as
a formal series. Since $B$ is entire, its formal differential identity
holds analytically by termwise differentiation near zero and then by
the identity theorem. Lemma~\ref{lem:finiteorder} gives a contradiction.
The equivalence between D-finite series and P-recursive coefficients
follows from the coefficient/operator translations in the proof of
Lemma~\ref{lem:factorialD}; hence the second conclusion.
\end{proof}

\begin{corollary}[Non-D-finiteness of feedback and its inverse]
\label{cor:nondfinite}
Under the hypotheses of Theorem~\ref{thm:borelgrowth}, neither $\U$
nor $\Q$ is D-finite. In particular, for $\lambda_j=j^2$ neither
series satisfies a nonzero linear differential equation with polynomial
coefficients, and neither coefficient sequence is P-recursive.
\end{corollary}
\begin{proof}
Apply Theorem~\ref{thm:obstruction} to the forward and inverse refined
types given by Theorem~\ref{thm:refinedfeedback}, with $d=p-\beta$.
\end{proof}

\begin{corollary}[No analytic change to a D-finite normal form]
\label{cor:normalform}
Let $F$ satisfy \eqref{eq:obstructionhyp}. Let $g$ be analytic with
$g(0)=0$, $g'(0)\ne0$, and let $H$ be analytic near $F(0)$ with
$H'(F(0))\ne0$. Then $H\circ F\circ g$ is not D-finite.
\end{corollary}
\begin{proof}
Apply Theorem~\ref{thm:coordinates} to the weight
$(n!)^m\log(n+\ee)^{-dn}$, with a finite log-convex modification as
in Lemma~\ref{lem:refinedweight}. The transformed type is
$|g'(0)|K$, which is positive and finite. Thus the transformed series
again satisfies the hypothesis of Theorem~\ref{thm:obstruction}.
\end{proof}

The last corollary is a growth obstruction, not a claimed general
closure property of the D-finite class under arbitrary analytic
composition. It says that the particular divergent growth profile
cannot be removed by invertible analytic coordinate changes.

\subsection{Limits of the obstruction}

Non-D-finiteness does \emph{not} mean differential transcendence. A
series can fail every linear polynomial-coefficient differential
equation and still satisfy a nonlinear algebraic differential equation.
No exclusion of such nonlinear equations is proved here.

The logarithmic deficit $d>0$ is also essential to this proof. At
$\beta=p$, the normalized Borel series has a finite positive radius
rather than being entire. For $\beta>p$, the same normalization has
zero radius. Lemma~\ref{lem:finiteorder} cannot be applied in either
case. Likewise, for noninteger $p-1$, the polynomial-recurrence
argument for dividing by $(n!)^{p-1}$ is unavailable as stated. These
are boundaries of the present theorem, not assertions that the series
at those boundaries are D-finite.

\section{Exact computations and reproducibility}
\label{sec:verification}

\subsection{Independent finite constructions}

The accompanying program \texttt{code/verify.py} uses only the Python
standard library. All identity and inequality checks use
\texttt{fractions.Fraction}; decimal diagnostics are generated
separately and labeled as floating point.

The forward coefficients are computed by a triangular exponential
recurrence. If
\[
 E_j(q)=\exp(\lambda_j\U(q))=\sum_{r\ge0}E_{j,r}q^r,
\]
then
\begin{equation}
 E_{j,0}=1,\qquad
 E_{j,r}=\frac{\lambda_j}{r}\sum_{k=1}^r k u_kE_{j,r-k},\qquad
 u_n=\sum_{j=1}^nE_{j,n-j}.
 \label{eq:verificationrecurrence}
\end{equation}
This construction is separate from the inverse construction, which uses
\eqref{eq:lagrange}. Both inverse compositions are subsequently checked
by direct polynomial multiplication. Equation~\eqref{eq:inversekernel}
is checked independently by substituting the computed inverse into the
literal finite truncation of the feedback kernel.

The finite majorant \eqref{eq:finitemajorant} is tested on signed rational
input series for four different log-convex weights: $n!$, $(n!)^2$,
$2^{n^2}$, and a slowly growing weight with nondecreasing quotients
$1+\operatorname{bitlength}(n)$. That last choice tests the mechanism
beyond factorial growth without using transcendental arithmetic.
The geometric-mean weight inequalities are checked after raising to
integer powers, avoiding numerical roots.

\subsection{Recorded results}

The recorded command was
\begin{quote}\ttfamily
python code/verify.py --order 70
\end{quote}
and it completed successfully. The exact report contains:
\begin{itemize}
\item Forward and inverse quadratic coefficients through degree $70$;
both inverse compositions, the literal inverse kernel, linear rescaling,
and the Catalan endpoint example verified through degree $30$.
\item Twelve seeded signed-series cases, with $180$ exact checks of the
finite inverse majorant and $8\,188$ exact concentration inequalities.
\item Geometric-mean weight checks for four weights through degree $12$,
and an analytic postcomposition/inverse factorization through degree $30$.
\end{itemize}
The random seed is $20260929$. The program raises an exception on a failed
check; it does not silently convert failed tests into warnings.
Table~\ref{tab:coefficients} gives the first ten exact coefficients.

\begin{table}[htbp]
\centering
\begin{tabular}{rll}
\toprule
$n$ & $u_n=[q^n]\U(q)$ & $q_n=[u^n]\Q(u)$\\
\midrule
1 & $1$ & $1$ \\
2 & $2$ & $-2$ \\
3 & $15/2$ & $1/2$ \\
4 & $107/3$ & $-2/3$ \\
5 & $1555/8$ & $-29/8$ \\
6 & $17362/15$ & $-743/60$ \\
7 & $5288311/720$ & $-8491/144$ \\
8 & $15414767/315$ & $-92726/315$ \\
9 & $2734575197/8064$ & $-611063/384$ \\
10 & $983674676/405$ & $-1667615261/181440$ \\
\bottomrule
\end{tabular}
\caption{Quadratic feedback $\lambda_j=j^2$: exact rational coefficients.}
\label{tab:coefficients}
\end{table}

These tests cannot verify a limsup over all indices or a statement about
entire growth at infinite radius. In particular, the floating-point
refined-root diagnostics at degree $70$ are approximately $7.59947$ for
$\U$ and $6.75803$ for $\Q$, not close to the asymptotic invariant $4$.
This slow preasymptotic behavior is recorded rather than hidden. No
monotone convergence claim or extrapolation of the constant is based on
that table. The asymptotic conclusions depend on the proofs and on the
two stated forward inputs, not on finite numerical proximity.

\subsection{Contents and build procedure}

The package contains the article source and PDF, the verification
program, exact coefficient data, numerical diagnostics, the recorded
verification report, a source-provenance file, and a Makefile. The TeX
source includes an embedded bibliography and the exact coefficient table
generated by the program. It compiles independently of the data directory
and has no network dependency.

To rebuild from the package directory, run
\begin{quote}\ttfamily
python code/verify.py --order 70\\
pdflatex -interaction=nonstopmode -halt-on-error article.tex\\
pdflatex -interaction=nonstopmode -halt-on-error article.tex\\
pdflatex -interaction=nonstopmode -halt-on-error article.tex
\end{quote}
The verification program does not modify ProveIt or contact external
services. The emitted \texttt{verification.json} distinguishes exact
checks from non-certified floating-point diagnostics. PDF layout checks
and build status are recorded separately in
\texttt{data/build\_validation.json}.

\section{Further research questions and a formalization program}
\label{sec:questions}

The following questions are proposed by this article. They are not all
claimed to be previously published open problems. Each specifies a
concrete next theorem or obstruction to investigate.

\subsection{Optimal subexponential cost of reversion}

The bound $\exp\{O(n/\sqrt{r_n})\}$ is sufficient for exact type but
comes from deliberately coarse composition counts. What is the smallest
universal overhead for a fixed weight and fixed input constants? Can
one replace the square-root loss by an expression involving the local
quotients $M_n/M_{n-1}$, and construct signed or nonnegative examples
that attain it? A useful first target is $M_n=(n!)^s$ with $0<s<1$,
where many small nonlinear blocks may still contribute a visible
subexponential factor.

\subsection{Full inverse asymptotics for quadratic feedback}

For $\lambda_j=j^2$, determine whether the normalized inverse roots in
\eqref{eq:quadratic} actually converge to four. A stronger target is
an asymptotic equivalent for $q_n$, including its sign and
subexponential amplitude. The small-order sign pattern in
Table~\ref{tab:coefficients} is not a proof of an eventual sign law.
The two-sided inverse argument used here deliberately avoids that issue;
a complete answer likely requires control of cancellation between
Lagrange blocks.

\subsection{Angular Borel growth and summation}

Determine the directional growth of $\B\Q(re^{i\theta})$ rather than
its maximum over the circle. Theorem~\ref{thm:borelgrowth} supplies an
exact global limsup scale, while \cite{negativeray} supplies a different,
directional summation result. Can these be joined into an angular
indicator theory, including the transition between the left and right
half-planes? This is also a route toward testing angular rather than
fine tubular summability.

\subsection{Critical Borel singularities}

For $\lambda_j\sim A(j\log j)^2$, the radius $1/(4A)$ is now fixed
for the inverse. Which boundary points are singular, what are their
local singularity types, and how do they depend on lower-order slope
terms? Positivity of the forward coefficients helps locate a forward
singularity, but cannot be transferred to a signed inverse without a
separate argument. A radius theorem is not yet a Stokes-constant theorem.

\subsection{Weights beyond log-convexity}

Theorem~\ref{thm:iff} gives a complete answer \emph{within} the
log-convex class. Find necessary and sufficient conditions for arbitrary
positive weights to have the same exact reversion invariant. A possible
approach is to compare a weight with a log-convex envelope up to
$\exp(o(n))$ factors. The difficult cases are oscillatory weights whose
deep dips occur precisely at indices populated by nonlinear
compositions.

\subsection{Two divergent composition arguments}

Analytic pre- and postcomposition are controlled exactly here. Determine
what can be said when both series are divergent. A naive equality in
terms of the two types alone cannot ignore cancellations, since a series
composed with its inverse is analytic. A useful replacement may be a
filtered composition estimate together with a noncancellation condition,
or a two-parameter profile that retains the position of large
coefficients.

\subsection{Multivariate and operator-valued inversion}

Extend the zero-loss theorem to formal maps in several variables,
where the derivative is an invertible matrix and homogeneous degrees
have many monomials. The natural question is whether an exact type
transforms by a directional norm of the inverse Jacobian rather than a
single scalar. A Banach-space formulation should distinguish constants
coming from multilinear operator norms from combinatorial losses that
are only subexponential in the degree.

\subsection{Hahn supports and accumulating action sets}

For a well-based exponential support not identified with the positive
integers, choose a meaningful replacement for degree and for $M_n$.
Can an action filtration support an analogue of the concentration
inequality? If infinitely many primitive actions lie below a finite
cut, the finite ordered-composition count used in
Theorem~\ref{thm:finite} is no longer available verbatim. A theorem
there must explicitly separate support summability from coefficient
majorization.

\subsection{Differential transcendence and the critical boundary}

Strengthen Corollary~\ref{cor:nondfinite} from a linear obstruction to
an obstruction against nonlinear algebraic differential equations, or
find a nonlinear equation that the feedback model does satisfy. In
parallel, determine non-D-finiteness at $\beta\ge p$, where the entire
normalization argument fails, and for noninteger $p-1$. These are
separate questions; none is settled merely by the terminology
``non-holonomic.''

\subsection{Summation compatible with refined inversion}

For logarithmic Gevrey deficits, identify a moment or acceleration
procedure that matches the full weight rather than just its factorial
order. Under what analytic continuation and growth hypotheses does its
sum commute with inversion and retain the exact type? Such a theorem
would have to add analytic hypotheses absent from the purely formal
result proved here, and must not infer summability from a zero type.

\subsection{A staged Lean formalization target}

A practical formalization can be divided into four theorem-sized stages.
First prove the finite log-convex concentration inequalities and the
geometric-mean ratio bound over positive real sequences. Next formalize
\eqref{eq:inversefinite} and \eqref{eq:finitemajorant} for coefficients
of an inverse formal power series, with the sum over positive
compositions made explicit. Third obtain the subexponential bound and
extended-real limsup theorem, including zero and infinite cases and the
Catalan necessity example. Finally connect the result to the particular
feedback construction and formally prove or import verified versions of
Repository inputs~\ref{input:forwardtype} and
\ref{input:forwardasymptotic}.

The general theorem is independent of those last two inputs, so it is a
useful standalone milestone. The finite exact checks in this package
supply examples and regression tests, not substitutes for those Lean
proofs.

\section{Conclusion}

The central result is a complete zero-loss criterion for log-convex
coefficient weights: universal exact type preservation under formal
reversion is equivalent to superexponential growth of the weight roots.
The proof provides a finite majorant rather than only an existence
statement, and yields exact analytic-coordinate covariance without
moderate-growth or summability assumptions.

Applied to the stated forward feedback estimates, this resolves the
repository's sharp inverse-type gap. It fixes the inverse Borel radius,
preserves logarithmically refined constants despite signed cancellation,
and yields a non-D-finiteness obstruction for an integer-order family,
including quadratic feedback. The next genuinely unresolved steps are
finer inverse asymptotics, angular Borel growth, and extensions beyond
log-convex degree weights; they are not concealed inside the proved
limsup statements.

\appendix
\section{A formal residue derivation of Lagrange reversion}
\label{app:lagrange}

For completeness, let $f(t)=at+O(t^2)$, $a\ne0$, and let $g=f^{-1}$.
For a formal Laurent series with a finite principal part, denote by
$\res_t$ the coefficient of $t^{-1}$. Formal substitution gives
\begin{equation}
 \res_z A(z)=\res_t A(f(t))f'(t).
 \label{eq:reschange}
\end{equation}
It is enough to check monomials $A(z)=z^k$. For $k\ne-1$ the right
side is the residue of $(f^{k+1})'/(k+1)$ and is zero. For $k=-1$,
write $f=at(1+h)$; then
$f'/f=t^{-1}+h'/(1+h)$ has residue one. At each coefficient only
finitely many relevant terms enter, justifying linear extension.

For $n\ge1$, apply \eqref{eq:reschange} to
$A(z)=g(z)z^{-n-1}$:
\[
 [z^n]g(z)=\res_t t f'(t)f(t)^{-n-1}.
\]
The residue of the derivative of $t f(t)^{-n}$ is zero. Since
\[
 (t f^{-n})'=f^{-n}-ntf'f^{-n-1},
\]
we obtain
\begin{equation}
 [z^n]g(z)=\frac1n\res_t f(t)^{-n}
 =\frac1n[t^{n-1}]\left(\frac{t}{f(t)}\right)^n.
 \label{eq:lagrangegeneral}
\end{equation}
This gives \eqref{eq:lagrange} at $a=1$ and also makes the scaling
factor $a^{-n}$ visible. The usual marker form
$V=t\Phi(V)$ follows by putting $f(v)=v/\Phi(v)$ when $\Phi(0)$
is a unit. For Proposition~\ref{prop:feedbackcoef}, work first over
$\C((q))$: the constant $\Phi(q,0)=q/(1-q)$ is a unit in this
Laurent-series field, so the argument applies verbatim to inversion in
the marker variable. Its displayed coefficient formula has $q$-order
at least $k$ in marker degree $k$, and therefore descends to the
coefficientwise specialization at marker value one in $\C[[q]]$.
Thus the calculation remains formal even when the resulting series
diverges everywhere except zero.

\section{Pinned source record and scope of novelty review}
\label{app:sources}

The snapshot is \commit, timestamped 29 September 2026,
20:14:16 UTC. Repository bibliography links are pinned to that commit.
The project README and these files under
\texttt{Analysis/Transseries/docs/series-and-transseries/} were the
principal sources:
\begin{quote}\small\ttfamily
Exponential\_Feedback\_Regularity\_Classification/article.tex\\
Exponential\_Feedback\_Regularity\_Classification/README.md\\
Negative\_Ray\_Summation\_Exponential\_Feedback/README.md.
\end{quote}
The first predecessor retains older internal snapshot references from
before the directory migration; those are not substituted for the
snapshot above. No upstream files were changed. Full links, retrieval
scope, and dependency boundaries are recorded in \texttt{PROVENANCE.json}.

The documented boundary resolved here is exact inverse type, not merely
Gevrey membership. This does not establish absence of every result from
all previous publications or concurrent notes. The classical ingredients
are Lagrange inversion, log-convex majorization, root tests, and D-finite
coefficient translations. The proposed contribution is their zero-loss
combination, the sharp weight threshold, and the feedback consequences.

\begin{thebibliography}{9}
\small
\addcontentsline{toc}{section}{References}

\bibitem{repo}
Vladimir Reshetnikov, \emph{ProveIt}, repository snapshot
\commit, 29 September 2026.
\href{https://github.com/VladimirReshetnikov/ProveIt/tree/0c973d8f5e7ef4550f4df1bf486d890037fd9f14/Analysis/Transseries}{Pinned transseries repository}.

\bibitem{reporeadme}
ProveIt contributors, \emph{Transseries}, project README, same snapshot.
\href{https://github.com/VladimirReshetnikov/ProveIt/blob/0c973d8f5e7ef4550f4df1bf486d890037fd9f14/Analysis/Transseries/README.md}{Pinned project README}.

\bibitem{feedback}
\emph{A sharp regularity classification for countable exponential-feedback
transseries}, AI-assisted research report prepared for Vladimir Reshetnikov,
29 September 2026, ProveIt research package. Conventional, non-Lean
results; precise input labels are recorded in Section~\ref{sec:feedback}.
\href{https://github.com/VladimirReshetnikov/ProveIt/blob/0c973d8f5e7ef4550f4df1bf486d890037fd9f14/Analysis/Transseries/docs/series-and-transseries/Exponential_Feedback_Regularity_Classification/article.tex}{Pinned article source}.
Companion scope statement:
\href{https://github.com/VladimirReshetnikov/ProveIt/blob/0c973d8f5e7ef4550f4df1bf486d890037fd9f14/Analysis/Transseries/docs/series-and-transseries/Exponential_Feedback_Regularity_Classification/README.md}{Pinned scope statement}.

\bibitem{negativeray}
\emph{Negative-Ray Summation of Countable Exponential-Feedback
Transseries}, AI-assisted research report prepared for Vladimir Reshetnikov,
29 September 2026, ProveIt research package; README consulted.
\href{https://github.com/VladimirReshetnikov/ProveIt/blob/0c973d8f5e7ef4550f4df1bf486d890037fd9f14/Analysis/Transseries/docs/series-and-transseries/Negative_Ray_Summation_Exponential_Feedback/README.md}{Pinned negative-ray README}.

\bibitem{gessel}
Ira M. Gessel, \emph{Lagrange inversion}, Journal of Combinatorial Theory,
Series A \textbf{144} (2016), 212--249.
\href{https://doi.org/10.1016/j.jcta.2016.06.018}{doi:10.1016/j.jcta.2016.06.018}.
\url{https://arxiv.org/abs/1609.05988}.

\bibitem{rainerschindl}
Armin Rainer and Gerhard Schindl, \emph{Composition in ultradifferentiable
classes}, Studia Mathematica \textbf{224} (2014), no.~2, 97--131.
\href{https://doi.org/10.4064/sm224-2-1}{doi:10.4064/sm224-2-1}.
\url{https://arxiv.org/abs/1210.5102}.

\bibitem{stability}
Armin Rainer and Gerhard Schindl, \emph{Equivalence of stability
properties for ultradifferentiable function classes}, Revista de la Real
Academia de Ciencias Exactas, F\'isicas y Naturales, Serie A,
Matem\'aticas \textbf{110} (2016), no.~1, 17--32.
\href{https://doi.org/10.1007/s13398-014-0216-0}{doi:10.1007/s13398-014-0216-0}.
\url{https://arxiv.org/abs/1407.6673}.

\bibitem{sauzin}
David Sauzin, \emph{Introduction to 1-summability and resurgence},
lecture notes, 2014, arXiv:1405.0356.
\url{https://arxiv.org/abs/1405.0356}.
\end{thebibliography}
\end{document}
